We will prove the conjecture below. Let us first see that it implies Theorem~\ref{thm-lef}. The argument here is similar to the proof of Theorem~\ref{thm-PP} in \cite{PapadakisPetrotou}. In fact, it is very natural to extend this theorem to the case of pseudo-manifolds $\Delta$ as in \cite{APP}. Recall that in Section~\ref{sec-mixed} we defined for any oriented pseudo-manifold $\Delta$ the algebra $\bar{H}(\Delta)$. This is equal to the algebra $H(\Delta)$ if $\Delta$ is a homology sphere.

\begin{coro} \label{cor-PPconj}
Let $h \in K[x_1,\ldots, x_N]_m$ for  some $0 \leq m \leq n/2$, and let $I, J$ have $n, n - 2m$ components, respectively. Then
\[ \partial_I W_\Delta(h^2 x_J) = (W_\Delta(h \sqrt{x_I x_J}))^2 \]
if $\sqrt{x_I x_J}$ exists, and is otherwise zero.
\end{coro}

\begin{proof}
Write $h$ as a linear combination of monomials, $h=\sum_L b_L x_L$. Then
\[ W_\Delta(h^2 x_J) = W_\Delta ( \sum_L b_L^2 x_L^2 x_J) = \sum_L b_L^2 W_\Delta (x_L^2 x_J).\]
Applying the derivative $\partial_I$ to this and using  Theorem~\ref{thm-PPconj}, we get
\[  \sum_L b_L^2 \partial_I W_\Delta (x_L^2 x_J) = \sum_L b_L^2 (W_\Delta (x_L \sqrt{x_I x_J}))^2 = (W_\Delta ( \sum_L b_L x_L \sqrt{x_I x_J}))^2\]
if the square root exists, and zero otherwise.
\end{proof}

The previous corollary shows why Theorem~\ref{thm-PPconj} is well suited for proving anisotropy theorems in characteristic $2$. The expression $W_\Delta(h \sqrt{x_I x_J})$ on the right hand side is the \Po pairing between $h$ and $\sqrt{x_I x_J}$. If $W_\Delta(h^2 x_J)$ on the left hand side is zero, then $h$ is orthogonal to $\sqrt{x_I x_J}$ for all $I$.

\begin{coro} \label{cor-reduce}
Let $\Delta$ be a pseudo-manifold of dimension $n-1$ over $K$. Consider $g\in K[x_1,\ldots, x_N]_m$ such that 
\[ l^p g^2 = 0 \quad \text{in $\bar{H}^{p+2m}(\Delta)$}\]
for some  $0\leq p \leq n-2m$. 
Let $q = \lfloor \frac{p}{2}\rfloor$. Then
\[ l^q g = 0 \quad \text{in $\bar{H}^{q+m}(\Delta)$.}\]
\end{coro}

\begin{proof}
If $l^p g^2 = 0$ in $\bar{H}^{p+2m}(\Delta)$, then for any integer vectors $I$ and $J$ with respectively $n$ and $n-2m-p$ components,
\[ \partial_I W_\Delta(l^p g^2 x_J) = 0. \]
If $p$ is even, then $p/2 = q$ and by Corollary~\ref{cor-PPconj},
\[ \partial_I W_\Delta(l^p g^2 x_J) = \partial_I W_\Delta((l^q g)^2 x_J) = (W_\Delta( l^q g \sqrt{x_I x_J}))^2 = 0 \]
if the square root exists, and so
\[ W_\Delta( l^q g \sqrt{x_I x_J}) = 0. \]
Since $I$ and $J$ can be chosen such that $x_I x_J = x_L^2$  for any of the monomials $x_L$ generating $\bar{H}^{n-q-m}(\Delta)$,  we conclude that $l^q g = 0$ in $\bar{H}^{q+m}(\Delta)$.

If $p$ is odd, then $l^p g^2 = (l^q g)^2 \cdot l = (l^q g)^2 \sum_{i=1}^N x_i$. Applying Corollary~\ref{cor-PPconj}, we have
\[ \partial_I W_\Delta(l^p g^2 x_J) = \sum_{i=1}^N  (W_\Delta(l^q g \sqrt{ x_i x_I x_J} ))^2 \]
if the square root exists. Let $L$ have $n-q-m$ components, and let $j$ be one of these components of $L$. Choose $I$ and $J$ such that $x_I x_J = x_L^2/x_j$. Then $\sqrt{ x_j x_I x_J}  = x_L$ and $\sqrt{ x_i x_I x_J}$ does not exist when $i \neq j$, so in this case
\[ \partial_I W_\Delta(l^p g^2 x_J)
 =   (W_\Delta(l^q g x_L))^2 = 0. \]
Hence,
\[ W_\Delta(l^q g x_L) = 0 \]
for any $x_L$. We conclude again that $l^q g = 0$ in $\bar{H}^{q+m}(\Delta)$.
\end{proof}

We now prove a generalization of Theorem~\ref{thm-lef} in the case of pseudo-manifolds.

\begin{theorem}
Let $\Delta$ be a pseudo-manifold of dimension $n-1$ over $K$. Then the quadratic form $Q_l$ on $\bar{H}^m(\Delta)$ is anisotropic  for any $m\leq n/2$.
\end{theorem}

\begin{proof}
Let $g \in \bar{H}^m(\Delta)$ be an isotropic element for $Q_l$,
\[ l^{n-2m} g^2 = 0.\]
Corollary~\ref{cor-reduce} implies that $l^q g = 0$, where $q = \lfloor \frac{n-2m}{2}\rfloor$. In particular, $l^q g^2 = 0$. Continuing this way we reduce the power of $l$ to zero and hence $g=0$.
\end{proof}
